\section{Local packets below an envelope}\label{sec:waves}

Throughout this Section the smooth positive pair \((A,\rho)\), and hence
\(G\) and \(L\), is fixed. The parameter \(k\) tends to infinity, and
\(h=h_k\) satisfies \eqref{eq:envelope-bounds}. Constants may depend on
the fixed pair, the constants in those bounds, and finitely many
construction orders. They are independent of \(k\) and of the packet
center. A lower threshold for \(k\) may depend on the same data.

The phase, transport amplitudes, and cutoff follow the complex WKB
construction of localized quasimodes; compare
\cite[Section~3]{DenckerSjostrandZworski2004}. Here the envelope varies
with the frequency. We therefore construct finite Taylor solutions and
track their dependence on its derivatives and Hessian margin, keeping
the support, derivative, and real-jet exponents fixed as the residual
accuracy increases.

We use complex bilinear extensions of real quadratic forms. When the
norm of a complex vector is used instead, it is the Hermitian norm.
In ordinary coordinate derivative estimates, a fixed finite atlas
with compactly contained subcharts is understood.

\subsection{The phase Hessian}

We first explain how the envelope test \(J_h\) enters the phase equation.
At a center \(z\), use \(G\)-normal coordinates and write
\(p=\nabla_Gh(z)\), \(Q=(1+|p|^2)^{1/2}\), and
\(H=\nabla_G^2h(z)\). For a unit vector \(v\perp p\), the complex
vector
\[
 P=p+iQv
\]
satisfies \(P\cdot P=-1\). A phase \(\phi\) with
\(\phi(z)=h(z)\) and \(\dd\phi(z)=P\) therefore satisfies the
constant term of the eikonal equation
\(G^{ij}\partial_i\phi\,\partial_j\phi=-1\). Its first-order terms
require \(SP=0\), where \(S=\partial^2\phi(z)\): the first metric
derivatives vanish in normal coordinates. A positive gap
\(H-\operatorname{Re}S\) then places \(\operatorname{Re}\phi\) below
\(h\) to quadratic order and supplies localization.

Write \(p=|p|e_1\) and \(b=|p|/Q\), choosing any unit
\(e_1\perp v\) when \(p=0\). For a symmetric matrix \(S=U+iD\),
the identity \(S(be_1+iv)=0\) forces
\[
 b^2U(e_1,e_1)+U(v,v)=0.
\]
Consequently, the corresponding test on the required gap is
\[
 b^2(H-U)(e_1,e_1)+(H-U)(v,v)
 =b^2H(e_1,e_1)+H(v,v)=J_h(v).
\]
This explains the positivity condition imposed on the envelope. The
following lemma constructs a compatible Hessian with a quantitative
gap, including when the real phase gradient vanishes.

\begin{lemma}[A dominated characteristic Hessian]\label{lem:phase-hessian}
Let \(e_1,v\) be orthonormal vectors in \(\R^4\), let \(0\le b<1\), and
let \(H\) be a real symmetric matrix such that
\[
 |H|\le k^{4\delta},\qquad
 J'=b^2H(e_1,e_1)+H(v,v)\ge k^{-2\delta}.
\]
For all sufficiently large \(k\), there are real symmetric matrices
\(U,D\) satisfying
\begin{equation}\label{eq:phase-hessian}
 (U+iD)(b e_1+i v)=0,\qquad
 H-U\ge k^{-3\delta}I,\qquad |U+iD|\le k^{20\delta}.
\end{equation}
The threshold can be chosen independently of \(H,e_1,v,b\) under these
hypotheses.
\end{lemma}

\begin{proof}
For \(b>0\), the first identity in \eqref{eq:phase-hessian} is equivalent
to
\[
 Dv=bUe_1,\qquad De_1=-Uv/b.
\]
The two prescribed columns are compatible with symmetry precisely when
\begin{equation}\label{eq:waves-column-compatibility}
 b^2U(e_1,e_1)+U(v,v)=0.
\end{equation}

First suppose that \(b\ge k^{-12\delta}\), and set
\[
 c_0=\frac{J'}{1+b^2},\qquad U=H-c_0I.
\]
Then \eqref{eq:waves-column-compatibility} holds. To give the extension
of \(D\) explicitly, write its blocks relative to
\(\R e_1\oplus\R v\oplus E\), where \(E=\{e_1,v\}^{\perp}\).
Denote the corresponding blocks of \(U\) by
\(u_{11},u_{12},u_{1E},u_{2E},\ldots\). Take
\[
 D=
 \begin{pmatrix}
 -u_{12}/b & b u_{11} & -u_{2E}/b\\
 b u_{11} & b u_{12} & b u_{1E}\\
 -u_{E2}/b & b u_{E1} & 0
 \end{pmatrix}.
\]
The identity \(u_{22}=-b^2u_{11}\) verifies both required columns.
Moreover,
\[
 H-U=c_0I\ge\tfrac12k^{-2\delta}I,\qquad
 |U|\le Ck^{4\delta},\qquad |D|\le Ck^{16\delta}.
\]
These imply \eqref{eq:phase-hessian} for large \(k\).

Now suppose that \(0\le b<k^{-12\delta}\). Put \(K=k^{14\delta}\) and
take
\[
 U=-K(I-vv^{\mathsf t})+b^2Kvv^{\mathsf t},\qquad
 D=-bK(e_1v^{\mathsf t}+ve_1^{\mathsf t}).
\]
The characteristic identity follows directly. These formulas also
give \(D=0\) and \(Uv=0\) when \(b=0\), so no division by a small \(b\)
is involved in this case.

For the domination inequality, split \(H-U-k^{-3\delta}I\) into
\(\R v\oplus v^\perp\). Its \(v^\perp\) block is at least
\(\tfrac12KI\) for large \(k\). Its scalar \(v\) block is at least
\[
 k^{-2\delta}-k^{-20\delta}-k^{-10\delta}-k^{-3\delta},
\]
because
\[
 H(v,v)\ge J'-b^2|H|,\qquad
 b^2|H|\le k^{-20\delta},\qquad b^2K\le k^{-10\delta}.
\]
The Schur complement subtracts at most
\(C|H|^2/K\le Ck^{-6\delta}\).
The positive \(k^{-2\delta}\) term dominates every subtraction.
The matrix is therefore positive definite for large \(k\).
Finally, \(|U|\le Ck^{14\delta}\) and \(|D|\le Ck^{2\delta}\),
proving the remaining bound.
\end{proof}

We will solve polynomial equations along a complex direction whose
bilinear square can be very small. Only its nonzero Hermitian norm
matters for this purpose.

\begin{lemma}[Polynomial directional inverses]\label{lem:polynomial-inverse}
Let \(\mathcal P_m\) be the space of complex homogeneous polynomials of
degree \(m\) on \(\R^4\), equipped with the coefficient norm in a fixed
monomial basis. Fix \(0<c\le C<\infty\).
If \(q\in\C^4\) satisfies \(c\le |q|\le C\), then
\[
 q\cdot\partial:\mathcal P_{m+1}\longrightarrow\mathcal P_m
\]
has a right inverse with norm bounded in terms of \(m,c,C\).
No lower bound for \(|q\cdot q|\) is required.
\end{lemma}

\begin{proof}
Choose an index \(i\) with \(|q_i|\ge c/2\), write
\(q\cdot\partial=q_i\partial_i+D'\), and let \(I_i\) denote integration
in \(X_i\) with zero integration constant. On \(\mathcal P_m\), set
\[
 \mathcal K=\frac{D'I_i}{q_i},\qquad
 R_{q,m}=\frac{I_i}{q_i}\sum_{r=0}^{m}(-\mathcal K)^r.
\]
Each application of \(\mathcal K\) lowers the degree in the variables
other than \(X_i\). Thus \(\mathcal K^{m+1}=0\) on \(\mathcal P_m\).
Since
\[
 (q\cdot\partial)\frac{I_i}{q_i}=I+\mathcal K,
\]
the finite geometric sum gives \((q\cdot\partial)R_{q,m}=I\).
The stated coefficient bound follows from the finite formula and the
bounds on \(q_i^{-1}\) and \(q\). The four possible choices of \(i\)
have the same uniform bound.
\end{proof}

\subsection{Construction and estimates for a packet family}

For a complex packet \(w\) centered at \(z\), let
\[
 \mathbf j_z w(x)
 =e^{-kh(x)}
   \left(w(x),\frac{\partial_\xi w(x)}{kQ_z}\right)\in\C^5,\qquad
 Q_z=\sqrt{1+|\nabla h(z)|^2}.
\]
Here \(\xi\) denotes the single normal coordinate system chosen at \(z\);
the denominator \(Q_z\) is kept fixed when \(x\) varies. A family
\(\mathcal W_z\) determines a real \(5\)-by-\(2|\mathcal W_z|\) matrix
\(\mathcal J_z(x)\) whose columns are the real and imaginary parts of
these jets. Its least spanning singular value is the square root of
the least eigenvalue of
\(\mathcal J_z(x)\mathcal J_z(x)^{\mathsf t}\).

\begin{proposition}[Packets with arbitrary finite accuracy]\label{prop:packets}
Fix a nonnegative integer \(J\) and a real number \(M>0\). Finite
construction orders can be chosen, before taking \(k\) large, so that
every center \(z\) with \(t(z)\ge1/4\) has a family \(\mathcal W_z\) of
four smooth complex packets with the following properties.

Each packet is supported in an ordinary coordinate ball about \(z\)
of radius at most \(Ck^{-1/3-100\delta}\). For every fixed
nonnegative integer \(j\),
\begin{equation}\label{eq:packet-derivatives}
 |D_x^j w(x)|\le C_j k^{2j}e^{kh(x)}.
\end{equation}
For the prescribed orders,
\begin{equation}\label{eq:packet-residual}
 |D_x^j(L+k^2)w(x)|\le C e^{kh(x)}k^{-M},
 \qquad 0\le j\le J.
\end{equation}
The center jets and the nearby jets satisfy
\begin{equation}\label{eq:jet-span}
 \sigma_{\min}(\mathcal J_z(z))\ge c k^{-3/4},\qquad
 \sigma_{\min}(\mathcal J_z(x))\ge k^{-1}
 \quad\text{if }\operatorname{dist}(x,z)\le k^{-10},
\end{equation}
for sufficiently large \(k\). Any fixed smooth reference distance
may be used here. The same conclusion holds if the latter radius
is multiplied by a fixed constant.

The exponents in the support, derivative, and jet bounds are
independent of \(J,M\). Constants and lower frequency thresholds may
depend on the chosen finite orders.
\end{proposition}

\begin{proof}
We divide the construction into the selection of jets, finite
polynomial equations, and localization estimates.

\paragraph{Normal coordinates and the linear jets.}
For each center choose a \(G\)-orthonormal frame and use normal
coordinates \(\xi\) in that frame. Put
\[
 p=\nabla h(z),\qquad Q=Q_z,\qquad H=\nabla^2h(z).
\]
In these coordinates the tangent and cotangent components agree at
the center. The envelope bounds, with a little slack in their
exponents, give
\begin{equation}\label{eq:waves-center-bounds}
 Q\le k^{3\delta},\qquad |H|\le k^{4\delta}
\end{equation}
for large \(k\).

If \(p\ne0\), set \(e_1=p/|p|\), and choose a favorable unit test \(e_2\)
from \eqref{eq:envelope-bounds}. If \(p=0\), first choose such an
\(e_2\), then choose any unit \(e_1\perp e_2\).
Complete to an orthonormal basis \(e_1,e_2,e_3,e_4\), and put
\[
 \eta_k=k^{-9\delta},\qquad
 v_2=e_2,\quad
 v_3=\frac{e_2+\eta_k e_3}{\sqrt{1+\eta_k^2}},\quad
 v_4=\frac{e_2+\eta_k e_4}{\sqrt{1+\eta_k^2}}.
\]
These three vectors span \(e_1^\perp\) with least spanning singular
value at least \(c\eta_k\).

Use the four pairs
\[
 (p',v)=(p,v_2),(p,v_3),(p,v_4),
 \bigl(p+Qk^{-3/4}e_1,v_2\bigr).
\]
For each pair define
\[
 Q'=\sqrt{1+|p'|^2},\qquad P=p'+iQ'v,\qquad b=|p'|/Q'.
\]
Then \(p'=|p'|e_1\), \(P\cdot P=-1\), and
\begin{equation}\label{eq:waves-imaginary-size}
 1\le Q'/Q\le1+k^{-3/4}.
\end{equation}
Changing \(e_2\) to a \(v_j\) changes its Hessian value by at most
\(Ck^{-5\delta}\). If \(a=|p|\), the function \(a^2/(1+a^2)\)
changes by at most \(2k^{-3/4}\) when \(a\) is increased by
\(Qk^{-3/4}\). Indeed its derivative at \(s\ge a\) is at most
\(2(1+a^2)^{-3/2}\). The additional error in the Hessian test is
at most \(Ck^{-3/4+4\delta}\). Both errors are negligible compared
with the original \(ck^{-\delta}\) margin. Consequently
\[
 b^2H(e_1,e_1)+H(v,v)\ge k^{-2\delta}
\]
for every selected pair. Apply Lemma~\ref{lem:phase-hessian} to obtain
\(U,D\). Multiplying its characteristic identity by \(Q'\) gives
\((U+iD)P=0\).

\paragraph{Frozen charts and scaling.}
The choice of frame may depend on \(k\), but each packet uses one
fixed chart when spatial derivatives are taken. More explicitly, if
\(E\) is its chosen frame, the chart map is
\[
 \Psi_{z,E}(\xi)=\exp_z\left(\sum_i\xi_iE_i\right).
\]
For each fixed derivative order \(q\), the derivatives of these maps
and their inverses are uniformly bounded over all centers and
orthonormal frames on sufficiently small common coordinate balls.
This follows from smoothness and compactness of the orthonormal
frame bundle of the fixed metric. In particular, the chain rule
gives
\[
 |D_\xi^q(h\circ\Psi_{z,E})|
 \le C_q(G)\sum_{\ell=1}^q\|h\|_{C^\ell}
 \le C'_q k^{\delta(q+1)}.
\]
There are no spatial derivatives of the rule used to select \(E\).

Set
\begin{equation}\label{eq:waves-scaling}
 r_0=k^{-100\delta},\qquad X=\xi/r_0,\qquad
 \kappa=kr_0Q,\qquad
 h_{\rm sc}(X)=\frac{h(\Psi_{z,E}(r_0X))-h(z)}{r_0Q}.
\end{equation}
For every fixed \(q\ge2\),
\begin{equation}\label{eq:waves-scaled-envelope}
 |D_X^q h_{\rm sc}|
 \le C_q Q^{-1}k^{\delta(101-99q)}.
\end{equation}
All these derivatives are bounded independently of \(k,z\). The
constant may depend on \(q\) and on high derivatives of the fixed
metric.

In normal coordinates write
\(L=g^{ij}(\xi)\partial_{\xi_i\xi_j}
+\beta^i(\xi)\partial_{\xi_i}\).
After multiplication by \(r_0^2\), the operator \(L+k^2\) becomes
\[
 \mathcal L+\kappa^2/Q^2,\qquad
 \mathcal L=g^{ij}(r_0X)\partial_{X_iX_j}
        +r_0\beta^i(r_0X)\partial_{X_i}.
\]
Its coefficients have uniformly bounded derivatives of every fixed
order on a fixed \(X\)-ball. The normal-coordinate identities are
\(g^{ij}(0)=\delta_{ij}\) and
\(\partial_{\xi_\ell}g^{ij}(0)=0\).

\paragraph{Finite eikonal and transport equations.}
Choose an integer \(T\ge3\), to be specified below. We construct a
complex polynomial \(\Phi\) with
\[
 \Phi(0)=0,\qquad
 \partial_X\Phi(0)=P/Q,\qquad
 \partial_X^2\Phi(0)=r_0(U+iD)/Q
\]
such that
\[
 E_\Phi(X)
 =g^{ij}(r_0X)\partial_i\Phi\,\partial_j\Phi+Q^{-2}
\]
vanishes to order \(T\) at zero. Here \emph{vanishes to order \(T\)}
means that all derivatives of order less than \(T\) vanish.
The constant equation follows from \(P\cdot P=-1\), and the linear
equations follow from \((U+iD)P=0\) and the vanishing first derivatives
of the metric.

At homogeneous degree \(m\ge2\), the as-yet unspecified homogeneous
term \(\Phi_{m+1}\) enters the equation through
\(2(P/Q)\cdot\partial\,\Phi_{m+1}\). All other terms have already
been determined. The Hermitian norm of \(P/Q\) lies between \(1\)
and \(3\) for large \(k\), by \eqref{eq:waves-imaginary-size}.
Lemma~\ref{lem:polynomial-inverse} therefore solves each equation
with a uniformly bounded right inverse. Solving through degree
\(T-1\) gives the desired polynomial. Its coefficients are bounded
for each fixed \(T\): the prescribed Hessian is bounded by
\(k^{-80\delta}\), the coefficient jets of \(\mathcal L\) are bounded,
and only finitely many bounded algebraic operations are used.

Define
\[
 \mathcal T a
 =2g^{ij}(r_0X)(\partial_i\Phi)(\partial_ja)
       +(\mathcal L\Phi)a.
\]
Fix an integer \(s\ge0\). Choose polynomials \(a_0,\ldots,a_s\) such
that
\[
 a_0(0)=1,\qquad a_j(0)=0\quad(j\ge1),
\]
and each of
\[
 \mathcal Ta_0,\qquad
 \mathcal Ta_j+\mathcal La_{j-1}\quad(1\le j\le s)
\]
vanishes to order \(T\) at zero. To see this directly, solve
\(\mathcal Ta=f\) by homogeneous degrees. The new amplitude term
of degree \(m+1\) enters the degree-\(m\) equation through
\(2(P/Q)\cdot\partial\); the remaining terms use only coefficients
already chosen. Lemma~\ref{lem:polynomial-inverse} again applies.
First construct the entire \(a_0\), then \(a_1\), and so on.
Thus the second derivative in \(\mathcal La_{j-1}\) causes no circular
dependence. Constants at zero are free, as stipulated above.
For fixed \(T,s\), all polynomial coefficients are bounded uniformly
in \(k,z\). No bound uniform as \(T,s\to\infty\) is asserted or needed.

Choose once a real smooth cutoff \(\omega\), equal to one on
\(\{|X|\le1\}\), supported in \(\{|X|<2\}\), and taking values in
\([0,1]\). Define
\begin{equation}\label{eq:packet-definition}
 \begin{split}
 a(X)&=\sum_{j=0}^s\kappa^{-j}a_j(X),\\
 w(\Psi_{z,E}(r_0X))
 &=e^{kh(z)}e^{\kappa\Phi(X)}a(X)\omega(k^{1/3}X),
 \end{split}
\end{equation}
and extend \(w\) by zero. For large \(k\), its support is compactly
contained in the normal chart, so this extension is smooth and has
the asserted support radius. On \(|X|\le2k^{-1/3}\),
\begin{equation}\label{eq:waves-amplitude}
 a(X)=1+O(k^{-1/3}),\qquad |a(X)|\ge\tfrac12,
\end{equation}
where the constant depends on \(T,s\). This follows from the chosen
constant terms and \(\kappa\ge1\).

\paragraph{Localization.}
The linear part of
\(\operatorname{Re}\Phi-h_{\rm sc}\) is either zero or
\(k^{-3/4}e_1\cdot X\). Lemma~\ref{lem:phase-hessian},
\eqref{eq:waves-center-bounds}, and
\eqref{eq:waves-scaled-envelope} imply
\begin{equation}\label{eq:waves-localization}
 \operatorname{Re}\Phi(X)-h_{\rm sc}(X)
 \le k^{-3/4}|X|-c k^{-106\delta}|X|^2+C|X|^3.
\end{equation}
The third derivative bounds used here are uniform for each fixed
construction order. On the packet support the cubic term can be
absorbed into half the quadratic deficit, since
\(k^{-1/3+106\delta}\to0\).
The maximum of the remaining linear term minus the quadratic term
is \(O(k^{-3/2+106\delta})\). Also
\[
 k^{1-100\delta}\le\kappa\le k.
\]
Consequently
\[
 |e^{kh(z)+\kappa\Phi(X)}|
 \le 2e^{kh(\Psi_{z,E}(r_0X))}
 \quad (|X|\le2k^{-1/3})
\]
for large \(k\).

On the smaller annulus
\(\tfrac14k^{-1/3}\le|X|\le2k^{-1/3}\), the negative exponent has
magnitude at least \(ck^{1/3-206\delta}\), whereas the positive linear
exponent is at most \(Ck^{-1/12}\). Since
\(1/3-206\delta>1/4\), a larger threshold gives
\begin{equation}\label{eq:waves-exponential-tail}
 |e^{kh(z)+\kappa\Phi(X)}|
 \le e^{-2k^{1/4}}e^{kh(\Psi_{z,E}(r_0X))}
 \quad\left(\tfrac14k^{-1/3}\le|X|\le2k^{-1/3}\right).
\end{equation}

\paragraph{Residual and derivative estimates.}
For the uncut expression, direct differentiation gives
\[
 \begin{split}
 e^{-\kappa\Phi}
 (\mathcal L+\kappa^2/Q^2)(e^{\kappa\Phi}a)
 ={}&\kappa^2E_\Phi a+\kappa\mathcal Ta_0\\
 &+\sum_{j=1}^s\kappa^{1-j}
       (\mathcal Ta_j+\mathcal La_{j-1})
       +\kappa^{-s}\mathcal La_s.
 \end{split}
\]
Taylor's theorem bounds a remainder vanishing to order \(T\), after
\(q\le J\le T\) derivatives in \(X\), by \(C|X|^{T-q}\).
On the support this is at most \(Ck^{-(T-q)/3}\).
Each ordinary derivative of such a remainder therefore has a cost
at most \(k^{1/3+100\delta}\), relative to its zeroth-order Taylor
bound. A derivative falling on the exponential, amplitude, or cutoff
has a cost at most a fixed constant times
\[
 r_0^{-1}(1+\kappa+k^{1/3})\le Ck^{1+3\delta}<Ck^2.
\]
Higher coordinate derivatives obey the same conservative \(k^2\)
cost per derivative, by the chain rule and the fixed normal-chart
bounds.

Undoing the operator scaling costs \(r_0^{-2}=k^{200\delta}\).
Thus, for \(0\le j\le J\), the uncut residual divided by \(e^{kh}\)
is bounded by
\begin{equation}\label{eq:waves-accuracy-budget}
 C k^{200\delta+2J}
 \left(k^{2-T/3}+k^{1-T/3}
       +k^{-(1-100\delta)s}\right).
\end{equation}
The commutators with the cutoff, and any prescribed derivatives of
them, have polynomial factors times \eqref{eq:waves-exponential-tail}.
They are smaller than every inverse power of \(k\).

For example, it suffices to choose integers \(T,s\) with
\[
 T>3(M+200\delta+2J+3),\qquad
 s>\frac{M+200\delta+2J+1}{1-100\delta}.
\]
These choices are made before the frequency threshold.
Increasing that threshold absorbs all their finite constants and
proves \eqref{eq:packet-residual}. The same derivative calculation
without the operator factor proves \eqref{eq:packet-derivatives} for
every fixed \(j\), with exponent \(2j\) independent of \(T,s\).

\paragraph{Real jet spanning.}
At the center, \(a(0)=1\) and the cutoff is constant, so
\[
 \mathbf j_z w(z)=\left(1,\frac{P}{Q}\right)+O(\kappa^{-1}).
\]
The error is bounded by \(Ck^{-1+100\delta}\).
Before this error, the real parts contain the two vectors
\[
 (1,(|p|/Q)e_1),\qquad
 (1,(|p|/Q+k^{-3/4})e_1).
\]
Their coefficients are bounded, and their least spanning singular
value on the value--\(e_1\) plane is at least \(ck^{-3/4}\).
The imaginary parts span the orthogonal \(e_1^\perp\) slope space
with least singular value at least \(ck^{-9\delta}\).
Together their least spanning singular value is at least
\(ck^{-3/4}\). Since
\(k^{-1+100\delta}=o(k^{-3/4})\), the actual center jets retain this
bound, with a smaller constant. The argument applies without
change at \(p=0\).

Finally, \eqref{eq:packet-derivatives} through order two and the
first envelope derivative bound give
\[
 |D_x(\mathbf j_z w(x))|\le Ck^3
\]
near the center. In differentiating this expression \(Q_z\) remains
fixed, while differentiating \(e^{-kh(x)}\) costs
\(k|D_xh|\le Ck^{1+2\delta}\).
A displacement of \(O(k^{-10})\) therefore changes the jet matrix
by \(O(k^{-7})\). Uniformly bounded changes between ordinary and
normal coordinates preserve these estimates.
This is negligible relative to \(k^{-3/4}\), and proves the second
inequality of \eqref{eq:jet-span}, after increasing the threshold.
\end{proof}

\subsection{Deterministic estimates at the local oscillation scale}

The next Lemma supplies estimates for an arbitrary polynomial-sized
collection of the preceding packets. It includes packets whose
centers lie outside the region in which signs will be measured.
Use the fixed chart \(\Omega\subset\{t>1/2\}\) and the open set
\(V\Subset\Omega\) from the preceding Sections. Points of \(\Omega\)
are identified with their coordinate vectors.

\begin{lemma}[Packets on a wavelength ball]\label{lem:scaled-packets}
Fix the finite construction orders in Proposition~\ref{prop:packets}
and a fixed finite number \(K_1\ge0\). Let \(\mathcal W\) be the collection of
all four packets from each of finitely many centers in
\(\{t\ge1/4\}\), with \(|\mathcal W|\le Ck^{K_1}\).
Suppose that for every \(x\in V\) some center is within
\(C_0k^{-10}\) of \(x\), where \(C_0\) is fixed. Set
\[
 R_x=(kQ_x)^{-1},\qquad
 S_x^2=\sum_{w\in\mathcal W}|w(x)|^2,\qquad
 w_x(Y)=w(x+R_xY),\quad |Y|\le2.
\]
The coordinate ball in this definition lies in \(\Omega\) for large
\(k\).

Consider the packets whose supports meet this closed ball. Their
centered scaled coordinates \(X_z\) are defined on the whole ball.
Partition them into
\[
 \mathcal N_x=\{w:|X_z(x)|<\tfrac12k^{-1/3}\},\qquad
 \mathcal T_x=\{w:|X_z(x)|\ge\tfrac12k^{-1/3}\}.
\]
For every fixed nonnegative integer \(j\), uniformly in \(x\in V\),
\begin{equation}\label{eq:scaled-packet-bounds}
 \begin{aligned}
 S_x&\ge c e^{kh(x)},\\
 \sup_{|Y|\le2}|D_Y^j w_x(Y)|
     &\le C_j|w(x)| &&(w\in\mathcal N_x),\\
 \sup_{|Y|\le2}|D_Y^j w_x(Y)|
     &\le C_j e^{-k^{1/4}}e^{kh(x)}
                                  &&(w\in\mathcal T_x).
 \end{aligned}
\end{equation}
Every packet in \(\mathcal N_x\) has cutoff identically one and
nonzero amplitude on the ball. Moreover,
\begin{equation}\label{eq:imaginary-gradient}
 \left|\operatorname{Im}
   \left(\frac{\dd_Yw_x}{w_x}\right)(0)\right|\ge c
 \qquad(w\in\mathcal N_x).
\end{equation}
In particular,
\begin{equation}\label{eq:waves-square-sums}
 \begin{aligned}
 \frac1{S_x^2}\sum_{w\in\mathcal W}
   \sup_{|Y|\le2}|D_Y^jw_x(Y)|^2&\le C_j,\\
 \frac1{S_x^2}\sum_{w\in\mathcal N_x}|w(x)|^2&=1-o(1).
 \end{aligned}
\end{equation}
The constants and thresholds may depend on the fixed construction
orders and on \(K_1,C_0\), but not on \(k,x\).
\end{lemma}

\begin{proof}
The bounds on the first two derivatives of \(h\) imply
\[
 |D_xQ|\le Ck^{4\delta}
\]
in the fixed chart. A packet centered at \(z\) that meets the ball
about \(x\) has
\[
 |z-x|\le Ck^{-1/3-100\delta}+C/k.
\]
For large \(k\), such a \(z\), and the whole ball, belong to a
compactly contained part of \(\Omega\). Normal coordinates centered
at \(z\) are defined throughout the ball. The preceding estimates
give
\[
 |Q_z-Q_x|
 \le Ck^{-1/3-96\delta}+Ck^{-1+4\delta}=o(1).
\]
In particular, \(Q_z/Q_x\) is bounded above and below independently
of the packet. The same comparison holds for \(Q\) throughout the
ball. Since \(|\dd h|\le CQ\) in the coordinate chart,
\begin{equation}\label{eq:waves-local-weight}
 k|h(x+R_xY)-h(x)|\le C,\qquad |Y|\le2.
\end{equation}

The nearest center has
\(|X_z(x)|=O(k^{-10+100\delta})\).
The center formula and its nearby-jet estimate in the proof of
Proposition~\ref{prop:packets} show that one of its packets has
\(|w(x)|\ge c e^{kh(x)}\). This proves the first line of
\eqref{eq:scaled-packet-bounds}.

Write \(\kappa_z=kr_0Q_z\). For every fixed \(q\ge1\), the chain
rule and normal-chart bounds give
\[
 |D_Y^qX_z|\le C_qr_0^{-1}R_x^q,\qquad
 \kappa_z|D_Y^qX_z|
 \le C_q\,\frac{Q_z}{Q_x}(kQ_x)^{1-q}\le C_q.
\]
For \(q=1\) there is also a lower bound: the matrix
\begin{equation}\label{eq:waves-scaled-jacobian}
 \kappa_z D_YX_z=(Q_z/Q_x)D_x\xi_z
\end{equation}
has inverse bounded independently of \(k,x,z\).
In particular,
\[
 |X_z(x+R_xY)-X_z(x)|
 \le Ck^{-1+100\delta}/Q_x=o(k^{-1/3}).
\]

Suppose first that \(w\in\mathcal T_x\).
At every point where its cutoff or any derivative of the cutoff can
contribute,
\[
 \tfrac14k^{-1/3}\le |X_z|\le2k^{-1/3}.
\]
Apply \eqref{eq:waves-exponential-tail}. Each fixed positive
\(Y\)-derivative of \(\kappa_z\Phi(X_z)\) is bounded: in a term
containing \(m\) derivatives of \(X_z\), the displayed derivative
bounds leave a factor \(\kappa_z^{1-m}\le1\).
Amplitude derivatives are bounded as well.
Derivatives of \(\omega(k^{1/3}X_z)\) are bounded because
\(k^{1/3}/\kappa_z\le k^{-2/3+100\delta}\le1\).
Thus differentiation introduces only fixed constants at each fixed
order. Together with \eqref{eq:waves-local-weight}, this proves
the tail bound in \eqref{eq:scaled-packet-bounds}.
This argument uses an absolute estimate and is valid even if
\(w(x)=0\).

For \(w\in\mathcal N_x\), the displacement estimate shows that
\(|X_z|<k^{-1/3}\) throughout the ball, so its cutoff is identically
one. By \eqref{eq:waves-amplitude}, the amplitude is bounded away
from zero. The bounds just proved for derivatives of
\(\kappa_z\Phi(X_z)\) and \(a(X_z)\) imply bounded derivatives of
\(w_x\) relative to \(|w_x|\).
Integrating the bounded real logarithmic derivative along straight
segments in the \(Y\)-ball gives
\[
 |w_x(Y)|\le C|w_x(0)|=C|w(x)|.
\]
This proves the second line of \eqref{eq:scaled-packet-bounds},
including \(j=0\).

For the lower oscillation estimate, the bounded coefficients of
\(\Phi\) give
\[
 \operatorname{Im}\partial_X\Phi(X_z(x))
   =(Q'/Q_z)v+O(k^{-1/3}).
\]
By \eqref{eq:waves-imaginary-size} its leading term has norm at
least one. The inverse bound in \eqref{eq:waves-scaled-jacobian}
therefore gives a positive lower bound for the imaginary
\(Y\)-gradient of \(\kappa_z\Phi(X_z)\).
The amplitude contributes only
\[
 \left|\frac{\dd_Ya(X_z)}{a(X_z)}\right|
 \le C r_0^{-1}R_x=Ck^{-1+100\delta}/Q_x=o(1).
\]
It cannot cancel that lower bound.
This proves \eqref{eq:imaginary-gradient}, including at centers
with \(p=0\) and for phases nearly isotropic in the complex
bilinear form.

Packets whose supports do not meet the ball contribute zero to
all the displayed sums. For the other packets, squaring
\eqref{eq:scaled-packet-bounds}, dividing by
\(S_x^2\ge c e^{2kh(x)}\), and using the cardinality bound shows
that the total tail contribution is at most
\[
 C_j k^{K_1}e^{-2k^{1/4}}=o(1).
\]
The non-tail contributions are bounded by a constant times
\(\sum_{\mathcal N_x}|w(x)|^2/S_x^2\le1\).
At derivative order zero the definition of \(S_x\) also shows that
this latter ratio equals \(1-o(1)\).
This proves \eqref{eq:waves-square-sums}.
\end{proof}
